% Generated draft: requires mathematical and source review.
\documentclass[11pt]{article}
\usepackage[margin=25mm]{geometry}
\usepackage{fontspec}
\setmainfont{Latin Modern Roman}
\usepackage{amsmath,amssymb,mathtools,unicode-math}
\setmathfont{Latin Modern Math}
\usepackage{enumitem,needspace,xurl,hyperref}
\setlength{\parindent}{0pt}
\setlength{\parskip}{5pt}
\setlength{\emergencystretch}{3em}
\newcommand{\N}{\mathbb{N}}
\newcommand{\Z}{\mathbb{Z}}
\newcommand{\Q}{\mathbb{Q}}
\newcommand{\R}{\mathbb{R}}
\newcommand{\C}{\mathbb{C}}
\newcommand{\F}{\mathbb{F}}
\newcommand{\A}{\mathbb{A}}
\newcommand{\PP}{\mathbb{P}}
\newcommand{\E}{\mathbb{E}}
\newcommand{\eps}{\varepsilon}
\newcommand{\dd}{\,\mathrm{d}}
\begin{document}
\section*{UnsolvedMath ID: 30003528}
\textbf{Source problem code:} catalogue entry
\subsection{Definitions and mathematical statement}
Let \(M\) be a closed smooth three-manifold, where closed means compact and without boundary, and let \(\lambda\) be a smooth contact form on \(M\), meaning that \(\lambda\wedge\dd\lambda\) is nowhere zero. Its Reeb vector field \(R_\lambda\) is uniquely characterized by
\[
\lambda(R_\lambda)=1,\qquad \dd\lambda(R_\lambda,\cdot)=0.
\]
A simple closed Reeb orbit is a periodic trajectory \(\gamma:\R/T\Z\to M\) satisfying \(\dot\gamma=R_\lambda\circ\gamma\), where \(T>0\) is its least positive period. Count such orbits by their geometric images, so that reparametrizations and multiple covers do not count as additional simple orbits. Define
\[
\mathcal P(M,\lambda)=\{\gamma(\R/T\Z):\gamma\text{ is a simple closed Reeb orbit}\},\qquad
N(M,\lambda)=\#\mathcal P(M,\lambda).
\]
\textbf{Source question.} Determine what restrictions hold on \(N(M,\lambda)\), and which finite or infinite cardinalities can occur, as \((M,\lambda)\) ranges over closed three-manifolds equipped with contact forms. No nondegeneracy or first Chern class assumption is imposed in the supplied question.

The supplied literature assessment identifies the two-or-infinity alternative
\[
N(M,\lambda)=2\quad\text{or}\quad \mathcal P(M,\lambda)\text{ is infinite}
\]
as a particular focus. The original question is broader and does not itself state this alternative as a conjecture. Moreover, it does not specify connectedness. The alternative must therefore not be adopted as a universal formulation without clarifying the intended connectedness convention: orbit counts add over disjoint components. Here ``infinite'' does not require countability.

For interpreting the supplied assessment, write \(\xi=\ker\lambda\). The form \(\dd\lambda|_\xi\) orients this plane bundle and defines its first Chern class \(c_1(\xi)\in H^2(M;\Z)\) using any compatible complex structure. This class is torsion if \(k c_1(\xi)=0\) for some positive integer \(k\). A contact form is nondegenerate if, for every periodic Reeb orbit, including multiple covers, the linearized return map on \(\xi\) has no eigenvalue equal to \(1\).

The supplied dated curation reports that the two-or-infinity alternative has been proved for nondegenerate contact forms and for all contact forms with torsion first Chern class, while the unrestricted degenerate case is not covered by those theorems. This is a supplied literature assessment, not an independently verified theorem statement; the exact hypotheses and connectedness conventions of the cited results require review.
\subsection{Short English statement}
Count the geometrically distinct simple periodic trajectories of the Reeb vector field of a contact form on a closed three-manifold. The source asks broadly what numbers can occur. Its supplied curation highlights the possibility of exactly two or infinitely many, but the intended connectedness convention and the precise hypotheses of the cited results require clarification. The supplied status is partially solved.
\subsection{Sources}
\begin{itemize}
\item Local UnsolvedMath registry, supplied complete source record: permanent ID \(30003528\), source problem code \texttt{catalogue entry}, title ``Number of Simple Reeb Orbits on Three-Manifolds,'' in the Oberwolfach Reports Open Problems collection. The recorded proposed year is 2017; no individual proposer is supplied.
\item Supplied workshop citation: \emph{Dynamische Systeme (2018)}, Oberwolfach Reports, pp.\ 1987--2046, DOI \texttt{10.4171/owr/2017/32}. \url{https://doi.org/10.4171/owr/2017/32}. The literature-source list labels this reference ``Oberwolfach Report 32/2017.'' The differing year labels are retained as supplied.
\item Colin--Dehornoy--Rechtman, \emph{Supporting broken books}, as labeled in the supplied literature sources. \url{https://arxiv.org/abs/2001.01448}.
\item Cristofaro-Gardiner et al., \emph{Two-or-infinity conjecture}, as labeled in the supplied literature sources. \url{https://arxiv.org/abs/2310.07636}.
\item Literature checked at: \texttt{2026-08-21}. This is supplied dated curation, not independently verified. It labels the problem partially solved and reports the nondegenerate and torsion-first-Chern-class cases described above, without claiming that those results cover the unrestricted degenerate case.
\item The registry labels the statement \texttt{corrected\_verified} and records a manual semantic, typographic, and LaTeX review. These are supplied provenance claims, not mathematical verification of this draft or independent verification of the literature.
\end{itemize}
\end{document}
